Cette thèse de doctorat aborde des enjeux fondamen-
taux à l’intersection de l’intelligence artificielle, de la
société et de l’éthique, en examinant quatre questions
structurantes.
Premièrement, elle explore le concept
d’**interprétabilité des modèles d’apprentissage
automatique**. Face à la complexité croissante des
algorithmes dits "boîtes noires", l’interprétabilité n’est
plus une option mais une nécessité, particulièrement
dans des domaines sensibles comme la santé ou la
justice. En nous basant sur les travaux de Doran et al.
(2017) et Burkart and Huber (2021), nous définissons
ce concept multidimensionnel, analysons ses enjeux
éthiques et pratiques, et disséquons les forces et
faiblesses des principales méthodes techniques (e.g.,
LIME, SHAP), soulignant les compromis inhérents
entre fidélité, intelligibilité et complexité.
Deuxièmement, la thèse se penche sur les **systèmes
de recommandation**, devenus omniprésents dans
l’écosystème numérique. À partir des études de Ko et al.
(2022) et Koren and Bell (2011), nous présentons une
taxonomie des techniques existantes avant de comparer
en profondeur deux approches paradigmatiques : le fil-
trage basé sur le contenu et le filtrage collaboratif par
factorisation de matrices. L’analyse met en lumière leur
complémentarité, le premier offrant l’explicabilité et
une solution au démarrage à froid des items, le second
une précision et une capacité de découverte supérieures,
au prix d’une opacité accrue et d’une sensibilité au
manque de données.
Troisièmement, nous abordons une question cognitive
fondamentale en établissant un parallèle critique entre
les **algorithmes de personnalisation et le phénomène
bureaucratique**. En mobilisant les théories de Max
Weber sur la rationalité et la "cage de fer" (Weber,
1922; Serpa and Ferreira, 2019) et la critique de Da-
vid Graeber sur la "stupidité structurelle" et la violence
symbolique (Graeber, 2015), nous soutenons que les
algorithmes ne sont pas une rupture mais une intensi-
fication, une privatisation et une automatisation de la
logique bureaucratique. La "cage de fer" wébérienne
se métamorphose en une "boîte noire" algorithmique,
posant une crise de redevabilité (*accountability*) et
reconfigurant les dynamiques de pouvoir et de contrôle
social.
Enfin, la quatrième partie propose des cadres d’action
pour instiller la **responsabilité dans la pratique**.
Face aux logiques prédatrices que ces systèmes peuvent
engendrer, nous esquissons une méthodologie d’inter-
vention éthique pour leur modification, ancrée dans les
principes de la Déclaration de Montréal pour une IA
responsable. Puis, inspirés par le processus collabo-
ratif du Biosphere Code (Bio, 2015), nous formulons
les principes et la structure d’un code de déontologie
pour les concepteurs, conçu comme un outil culturel
pour naviguer les rapports de force organisationnels et
cultiver une responsabilité partagée, au-delà de la seule
imputabilité individuelle (Neuer, 2018).
En conclusion, cette thèse soutient que les défis po-
sés par l’IA moderne sont de nature socio-technique et
exigent des réponses qui transcendent la seule ingénierie. 
Elle appelle à une gouvernance plus démocratique
et éthique de ces nouvelles formes de pouvoir algorith-
mique.
Mots-clés : Interprétabilité de l’IA (XAI), Systèmes de
Recommandation, Filtrage Collaboratif, Bureaucratie
Algorithmique, Max Weber, David Graeber, Éthique de
l’IA, Responsabilité Algorithmique, Code de Déontolo-
gie.


Première partie
Introduction et Contexte
Général
Introduction Générale
1 La Révolution Silencieuse de
l’Apprentissage Automatique
L’aube du XXIe siècle est marquée par une transformation
technologique aussi profonde que silencieuse : l’avènement
de l’apprentissage automatique (AA) comme force motrice
de l’innovation et de l’organisation sociale. Des diagnostics
médicaux assistés par ordinateur à la gestion des chaînes lo-
gistiques mondiales, de la conduite de véhicules autonomes
à la curation personnalisée de nos flux d’information, les
algorithmes d’AA ne sont plus des curiosités de laboratoire.
Ils constituent une nouvelle infrastructure invisible qui mé-
diatise une part croissante de nos interactions avec le monde,
nos institutions et autrui.
Cette révolution repose sur une idée d’une simplicité trom-
peuse : la capacité des machines à apprendre des patterns
à partir de données, sans être explicitement programmées
pour chaque tâche spécifique. Cette capacité a débloqué des
performances prédictives et décisionnelles sans précédent,
en particulier avec l’émergence des modèles d’apprentissage
profond (deep learning). Cependant, cette puissance s’est
acquise au prix d’une complexité exponentielle. Les architec-
tures de réseaux de neurones profonds, avec leurs millions,
voire milliards de paramètres, fonctionnent souvent comme
des boîtes noires (black boxes) : leurs mécanismes internes,
qui transforment une entrée (par exemple, une image) en
une sortie (par exemple, un diagnostic), sont si complexes
qu’ils défient l’intuition et l’analyse humaine directe.
2 La Crise de la Boîte Noire : Confiance,
Responsabilité et Pouvoir
Cette opacité n’est pas un simple défi technique; elle
engendre une véritable crise conceptuelle et sociétale. Com-
ment pouvons-nous faire confiance à un système dont nous
ne comprenons pas le raisonnement ? Comment attribuer la
responsabilité en cas d’erreur ou de décision préjudiciable ?
Comment s’assurer que ces systèmes, déployés à grande
échelle, ne perpétuent pas ou n’amplifient pas les biais et les
injustices existants dans nos sociétés ?
Ces questions sont au cœur de cette thèse. Elles nous
obligent à dépasser une vision purement instrumentaliste
de la technologie pour l’examiner comme un phénomène
social et politique, porteur de nouvelles logiques de pouvoir
et de contrôle. La "boîte noire" n’est pas seulement opaque,
elle est puissante. Elle classifie, elle trie, elle recommande,
elle décide et, qui plus est, l’un des enjeux cen-
traux de notre époque est qu’il devient « impossible pour un
individu de mettre en doute une décision automatique prise
